% ==============================================================================
% CAPÍTULO 5 - CONSIDERAÇÕES FINAIS
% ==============================================================================

\chapter{Considerações Finais}
\label{cap:consideracoes}

Neste capítulo, você deve encerrar o relatório fazendo um balanço crítico da sua experiência. 
Escreva um parágrafo inicial retomando o objetivo geral do estágio e confirme se ele foi alcançado.
É o momento de refletir sobre como as atividades desenvolvidas contribuíram para o seu crescimento 
profissional e para a solução dos problemas da empresa.

%-------------------------------------------------
% Conclusões
%-------------------------------------------------
\section{Conclusões}
\label{sec:conclusoes}

Descreva os principais aprendizados e resultados obtidos. Destaque como a teoria da universidade
se conectou com a prática do mercado e quais foram as lições mais valiosas (técnicas ou comportamentais). 
Evite repetir descrições de tarefas; foque na análise do valor que o estágio agregou à sua formação e à instituição concedente.


%-------------------------------------------------
% Trabalhos Futuros
%-------------------------------------------------
\section{Trabalhos Futuros}
\label{sec:trabalhos-futuros}

Sugira caminhos para a continuidade do trabalho, seja em novos projetos na empresa, 
melhorias no sistema desenvolvido ou novas linhas de pesquisa acadêmica. 
Seja específico sobre o que poderia ser feito a seguir e como isso agregaria valor.

\begin{itemize}
    \item Sugestão de trabalho futuro 1;
    \item Sugestão de trabalho futuro 2;
    \item Sugestão de trabalho futuro 3.
\end{itemize}
